% ============================================================
%  PLANTILLA DE CLASE — ESUMER / Estud-IA
%  Curso: Análisis y Visualización de Datos para la Toma de Decisiones
%
%  Cómo usarla:
%   1. Copiar este archivo a algo como  clases/semana01-clase01.tex
%   2. Completar los campos de la sección "DATOS DE LA CLASE"
%   3. Escribir el contenido en las diapositivas de ejemplo
%      (Objetivos, Contenido, Actividad/Producto, Cierre)
%   4. Compilar con:  xelatex clases/semana01-clase01.tex
%      (requiere xelatex o lualatex por el uso de fontspec)
% ============================================================
\documentclass[aspectratio=169,11pt]{beamer}

% La ruta asume que se compila desde la raíz del curso (20262/).
% Si se compila desde otra carpeta, ajustar \graphicspath y el
% \usepackage{...} de la plantilla en consecuencia.
\usepackage{beamerthemeesumer}
\usetheme{esumer}

\usepackage[spanish]{babel}
\usepackage{booktabs}
\usepackage{graphicx}

% ---------------- DATOS DE LA CLASE ----------------
\title{Nombre del tema de la clase}
\subtitle{Subtítulo o gancho breve (opcional)}
\author{Docente: Nombre Apellido}
\date{Fecha de la sesión}

\setmodulo{Módulo 1: Explorar (Fundamentos y Consulta de Datos)}
\setsemana{1}
\setclase{1}
\setfechaclase{Sep 1 al 7}
\setmodalidad{Híbrido}
% -----------------------------------------------------

\begin{document}

{
\setbeamertemplate{footline}{}
\begin{frame}[plain]
  \titlepage
\end{frame}
}

\begin{frame}{Agenda}
  \begin{enumerate}
    \item Objetivos de la sesión
    \item Contenido / conceptos clave
    \item Actividad práctica
    \item Producto esperado
    \item Cierre y próxima sesión
  \end{enumerate}
\end{frame}

\section{Objetivos}
\begin{frame}{Objetivos de la sesión}
  \begin{block}{Al finalizar esta clase el estudiante podrá...}
    \begin{itemize}
      \item Objetivo 1.
      \item Objetivo 2.
      \item Objetivo 3.
    \end{itemize}
  \end{block}
\end{frame}

\section{Contenido}
\begin{frame}{Conceptos clave}
  \begin{itemize}
    \item Concepto 1.
    \item Concepto 2.
    \item Concepto 3.
  \end{itemize}
\end{frame}

\begin{frame}{Ejemplo / demostración}
  Espacio para explicación guiada, código o gráfico de ejemplo.
\end{frame}

\section{Actividad}
\begin{frame}{Actividad práctica}
  \begin{alertblock}{Reto de la sesión}
    Descripción de la actividad / taller / laboratorio a desarrollar.
  \end{alertblock}
\end{frame}

\begin{frame}{Producto esperado}
  \begin{exampleblock}{Entregable}
    Descripción del producto o evidencia que el estudiante debe entregar.
  \end{exampleblock}
\end{frame}

\section{Cierre}
\begin{frame}{Cierre y próxima sesión}
  \begin{itemize}
    \item Resumen de lo visto hoy.
    \item Tarea / preparación para la próxima clase.
    \item Recursos adicionales.
  \end{itemize}
\end{frame}

{
\setbeamertemplate{footline}{}
\begin{frame}[plain]
  \begin{tikzpicture}[remember picture,overlay]
    \fill[esumerNavy] (current page.south west) rectangle (current page.north east);
    \node at (current page.center) {\color{white}\Huge\bfseries ¡Gracias!};
  \end{tikzpicture}
\end{frame}
}

\end{document}
